\begin{frame}[noframenumbering,plain]
    \setcounter{framenumber}{1}
    \maketitle
\end{frame}

\begin{frame}
	\frametitle{Положения, выносимые на защиту}
	% \footnotesize
	\begin{itemize}
		\item Впервые показано, что ценовое поведение онлайн-ритейлеров РФ соответствует гибридным моделям ценообразования (state- и time-dependent).
		
		\item Установлена зависимость механизма ценообразования от инфляционного режима: при высокой инфляции преобладают state-dependent модели, при низкой — time-dependent.
		
		\item Выявлен инфляционный порог ($\approx$1\% в месяц), выше которого резко возрастает роль частоты пересмотров цен, что подтверждает необходимость её учёта в моделях.
		
		\item Жёсткость цен зависит от характеристик товаров и календарных факторов: выявлено отрицательное влияние торговой надбавки, различия по сроку хранения и сезонные колебания.
		
		\item Установлены технические и методологические ограничения сбора онлайн-данных о ценах, влияющие на интерпретацию эмпирических результатов.
	\end{itemize}
\end{frame}

\note{
    Проговариваются вслух положения, выносимые на защиту
}

\begin{frame}
    \frametitle{Содержание}
    \tableofcontents
\end{frame}
\note{
    Работа состоит из четырёх глав.

    \medskip
    В первой главе \dots

    Во второй главе \dots

    Третья глава посвящена \dots

    В четвёртой главе \dots
}
